\section{Uniform certificate tools}
\label{p14:cert:section}

Fixed-order estimates leave a separate problem: summing orders that grow
with $N$. The arithmetic certificates below solve it by counting
congruences and graph components with explicit cardinality weights.
Their definitions concern the actual values in the windows and apply
before any affine prescription is imposed. In particular, the geometric
counts also permit irregular windows; replacing $p$ by $p_x\le2p$ only
changes a fixed weight. The final parameter choices for the optimal
reconstruction are made in Section~\ref{p13:transfer:section}.
The complete elementary proofs are collected in
Section~\ref{p14:comp:certificate-tools}.

\subsection{Small kernels and moderate congruence certificates}

For $Y\ge B$, write
\[
 K_Y(v)=\prod_{\substack{\ell>Y\ \mathrm{prime}\\v_\ell(v)\ \mathrm{odd}}}\ell.
\]
The following elementary form of the Rankin argument will be used with
its exponent fixed.

\begin{lemma}[Counting a small truncated squarefree kernel]
\label{p1:core:rankin}
For fixed $1/2<\sigma<1$, fixed $C_1>0$, and $T\ge1$,
\begin{equation}
 \#\{v\le C_1N:K_Y(v)\le T\}
 \ll_{\sigma,C_1}N^\sigma T^{1-\sigma}
       \exp(C_\sigma Y^{1-\sigma}).
 \label{p1:core:rankin-bound}
\end{equation}
\end{lemma}
\noindent The complete proof is given in Section~\ref{p14:comp:certificate-rankin}.


Fix
\begin{equation}
 A_0>1,\quad 0<\beta<1,\quad \tfrac12<\sigma<1,
 \quad A_0(1-\sigma)<1,
 \label{p1:core:parameter-conditions}
\end{equation}
and put $H=\log N$, $Y=H^{A_0}$, $T=N^\beta$ and
$\delta=(1-\sigma)(1-\beta)>0$. Mark a window if one of its values
has $K_Y(v)\le T$. Each value lies in at most $B$ windows, so the
number $M$ of marked centres satisfies
\begin{equation}
 M\le N^{1-\delta+o(1)}.
 \label{p1:core:marks}
\end{equation}
Here and below the $o(1)$ depends on the fixed parameters, but not on
the size of a family subsequently summed over.

Choose two distinct values in an unmarked anchor window. Their kernels
are coprime, because a prime exceeding $Y\ge B$ cannot divide their
nonzero difference. Thus their product $Q$ is squarefree and
$T^2<Q\le9N^2$. If a relation selects these values, each prime of $Q$
occurs to odd order in at least one other window. Assign it to one such
occurrence. For each receiving block $j$, let $D_j$ be the product of
the assigned primes. The $D_j$ are coprime and $\prod_jD_j=Q$.
Call the resulting certificate \emph{moderate} when every $D_j\le6NB$.
After specifying a certificate we may discard the relation itself:
the congruences alone bound the possible centres.

\begin{lemma}[Counting moderate certificates]
\label{p1:core:moderate}
Let $\mathcal U_r$ be the ordered lists of $r$ distinct original windows having
such a certificate, and set
$m=\lceil 2\log(3N)/\log Y\rceil$. Then
\begin{equation}
 |\mathcal U_r|
 \le36rB^4 2^{r-1}N^rT^{-2}(rB)^m.
 \label{p1:core:moderate-count}
\end{equation}
For every fixed $u>0$,
\begin{equation}
 \sum_{r\ge1}\frac{(up)^r}{r!}|\mathcal U_r|
 \le N^{-2\beta+4/A_0+o(1)}.
 \label{p1:core:moderate-baseline}
\end{equation}
\end{lemma}
\noindent The complete proof is given in Section~\ref{p14:comp:certificate-moderate}.


\subsection{Saturation has small degree in both directions}

Put a directed edge $x\to y$ when two values in window $x$ have a
kernel product possessing a divisor $D>6NB$ all of whose primes divide
values in window $y$. We impose no direct-sum condition on $S_x,S_y$.
Let $\Gamma$ be the underlying simple undirected graph.

\begin{lemma}[Incoming and outgoing saturation degrees]
\label{p1:core:degrees}
Uniformly in the fixed vertex,
\begin{equation}
 \deg^+(x),\deg^-(x)\le N^{2/A_0+o(1)},
 \qquad \Delta(\Gamma)\le N^{2/A_0+o(1)}.
 \label{p1:core:degree-bound}
\end{equation}
\end{lemma}
\noindent The complete proof is given in Section~\ref{p14:comp:certificate-degrees}.


If a primitive family has no moderate certificate, every unmarked
vertex has a neighbour in its induced saturation graph. To see this,
choose a nonzero relation, select two of its values at that vertex and
assign the primes to exterior occurrences. Some receiver must have
$D_j>6NB$. The choices made at different vertices need not agree.
Their existence in one deterministic graph is sufficient.

\begin{lemma}[A generating function with marked isolated vertices]
\label{p1:core:generator}
Let a simple graph have $n$ vertices, maximum degree $\Delta$ and
$M$ marked vertices. If $\mathcal F$ consists of the vertex sets
whose induced isolated vertices are all marked, then, for $z\ge0$,
\begin{equation}
 \sum_{D\in\mathcal F}z^{|D|}
 \le\exp\left\{Mz+nz\bigl((1+z)^\Delta-1\bigr)\right\}.
 \label{p1:core:generator-bound}
\end{equation}
More generally, for nonnegative vertex weights $w_v$ and marked set
$\mathcal M$,
\[
 \sum_{D\in\mathcal F}\prod_{v\in D}w_v
 \le \exp\left\{\sum_{v\in\mathcal M}w_v+
 \sum_vw_v\left(\prod_{u\sim v}(1+w_u)-1\right)\right\}.
\]
\end{lemma}
\noindent The complete proof is given in Section~\ref{p14:comp:certificate-generator}.

\subsection{Proofs of the uniform certificate tools}
\label{p14:comp:certificate-tools}

All definitions and parameter conditions are those of
Section~\ref{p14:cert:section}. The counts are geometric, and all
constants are uniform in the size of the family being summed.

\subsubsection{Proof of Lemma~\ref{p1:core:rankin}}
\label{p14:comp:certificate-rankin}

\begin{proof}
Write uniquely $v=k s a^2$, where $k=K_Y(v)$ and $s$ is the
squarefree product of the primes at most $Y$ occurring to odd order.
The inequality $1\le (C_1N/(ksa^2))^\sigma$ on the range being
counted gives an upper bound
\[
 (C_1N)^\sigma
 \left(\sum_{k\le T}k^{-\sigma}\right)
 \left(\sum_{a\ge1}a^{-2\sigma}\right)
 \prod_{\ell\le Y}(1+\ell^{-\sigma}).
\]
Enlarging the sum over $k$ to all positive integers only increases it.
The three factors are respectively $O_\sigma(T^{1-\sigma})$,
$\zeta(2\sigma)$, and at most
$\exp(\sum_{j\le Y}j^{-\sigma})$. This proves the bound.
\end{proof}

\subsubsection{Proof of Lemma~\ref{p1:core:moderate}}
\label{p14:comp:certificate-moderate}

\begin{proof}
There are at most $rNB^2$ choices for the anchor and its two offsets.
The product $Q$ has at most $m$ prime factors, and each has at most
$rB$ destinations. Once these assignments are fixed, the Chinese
remainder theorem leaves at most $1+N/D_j$ possibilities for centre
$j$. At most two of the $D_j$ exceed $N$, since $Q\le9N^2$.
For those two modules, replacing $1+N/D_j$ by $2N/D_j$ costs at most
$(6B)^2$. Therefore
\[
 \prod_{j\ne\mathrm{anchor}}(1+N/D_j)
 \le 2^{r-1}(6B)^2N^{r-1}T^{-2}.
\]
This proves \eqref{p1:core:moderate-count}; allowing collisions and
incompatible assignments only enlarges the count.

After multiplication by $(up)^r/r!$, the remaining series is, up to
a fixed factor,
\[
 N^{-2\beta}B^{m+4}\sum_{r\ge0}\frac{z^r r^{m+1}}{r!},
 \qquad z=2uNp=O_u(1).
\]
Since $r^{m+1}\le(m+1)!e^r$, this series is at most
$N^{-2\beta}B^{m+4}(m+1)!\exp(ez)$. Finally
\[
 m=\left(\frac2{A_0}+o(1)\right)\frac{H}{\log H},\qquad
 \log(B^m(m+1)!)=\left(\frac4{A_0}+o(1)\right)H.
\]
This yields \eqref{p1:core:moderate-baseline}. The factorial here
divides the orderings of a set of distinct windows; it is not an
assumption on the growth of a fixed-order arithmetic constant.
\end{proof}

\subsubsection{Proof of Lemma~\ref{p1:core:degrees}}
\label{p14:comp:certificate-degrees}

\begin{proof}
For an outgoing edge, choose two anchor offsets in $B^2$ ways. Each of
the at most $m$ primes is either unused or assigned to one of $B$
offsets in the receiver. A qualifying code has modulus $D>N$ and
therefore determines at most one receiving centre in $I_N$. This gives
$B^2(B+1)^m=N^{2/A_0+o(1)}$ possibilities.

For the incoming degree, fix the receiving window $y$ and let
$\Omega_y$ be the primes greater than $Y$ dividing any of its values.
Writing $R=\lceil B\log(3N)/\log Y\rceil$, we have $|\Omega_y|\le R$.
A code chooses two anchor offsets, at most $m$ primes from $\Omega_y$,
and one of the two offsets for each chosen prime. For a qualifying
code their product exceeds $N$, so its congruences determine at most
one anchor centre. Thus the number of possible incoming neighbours is
at most
\[
 B^2\sum_{j=0}^m\binom Rj2^j
 \le B^2(m+1)(2eR/m)^m
 =N^{2/A_0+o(1)}.
\]
The displayed binomial bound holds for large $N$, since $m\le R/2$.
Nonqualifying codes may be included in the number of codes; the
one-centre conclusion is applied only to qualifying ones. No receiving
offset needs to be chosen in this incoming count: the receiver has
already been fixed. Symmetrisation adds the two degree bounds.
\end{proof}

\subsubsection{Proof of Lemma~\ref{p1:core:generator}}
\label{p14:comp:certificate-generator}

\begin{proof}
Fix an admissible set $D$. Its isolated vertices are marked singletons.
On the other vertices choose an inclusion-minimal edge cover of the
induced graph. Every edge of this cover has an endpoint of degree one
in the cover, since otherwise that edge could be removed. Its
components are therefore nontrivial stars. Thus $D$ has a partition
into marked singletons and stars, each with one centre and a nonempty
subset of its neighbours as leaves.

Write $w(A)=\prod_{v\in A}w_v$. Choose one such partition for each $D$.
The product over all possible blocks of $1+w(A)$ contains a distinct
term of weight $w(D)$ for this chosen partition: its union recovers
$D$. Dropping disjointness only adds nonnegative terms. Since
$1+u\le e^u$ for $u\ge0$, the weighted sum is at most the exponential
of the total block weight. The singleton contribution is
$\sum_{v\in\mathcal M}w_v$. For a fixed star centre $v$, its leaf sum
is exactly
\[
 w_v\sum_{\varnothing\ne A\subseteq N(v)}\prod_{u\in A}w_u
   =w_v\left(\prod_{u\sim v}(1+w_u)-1\right).
\]
Counting the same two-vertex star from both centres is harmless.
This proves the heterogeneous assertion. Set $w_v=z$ and use
$\deg(v)\le\Delta$ to obtain \eqref{p1:core:generator-bound}.
There is no small-activity hypothesis. When $\Delta z<2$, the
elementary bound $e^u-1\le2u/(2-u)$ also gives the convenient
majorant $\exp\{Mz+2n\Delta z^2/(2-\Delta z)\}$.
\end{proof}
